% Shared pgfplots style for the Isosyll paper figures (paper-data). Okabe--Ito colour-blind-safe palette;
% every series also differs by marker shape so the figures survive greyscale printing.
% Inputs CSVs from \paperdatadir (default data/, relative to docs/paper/paper.tex).
\providecommand{\paperdatadir}{data/}%
\definecolor{okVermillion}{RGB}{213,94,0}%
\definecolor{okBlue}{HTML}{D9480F}% house accent: Isosyll is orange, never the secondary blue%
\definecolor{okGreen}{RGB}{0,158,115}%
\definecolor{okOrange}{RGB}{230,159,0}%
\definecolor{okSky}{RGB}{86,180,233}%
\definecolor{okPurple}{RGB}{204,121,167}%
\definecolor{okGrey}{RGB}{110,110,110}%
\definecolor{okInk}{RGB}{64,64,64}% neutral marks beside the accent okBlue (darker, so they differ in greyscale too)
\pgfplotsset{%
  ultra/base/.style={%
    font=\footnotesize, tick label style={font=\scriptsize}, label style={font=\footnotesize},%
    legend style={font=\scriptsize, draw=none, fill=none}, axis line style={okGrey},%
    tick style={okGrey}, major grid style={gray!25}, grid=major, clip=false,%
  },%
  ultra/rzero/.style={color=okVermillion, mark=square*, mark size=1.6pt, thick},%
  ultra/dense/.style={color=okBlue, mark=*, mark size=1.7pt, thick},%
  ultra/mr/.style={color=okGreen, mark=triangle*, mark size=2.1pt, thick},%
  ultra/failed/.style={mark options={fill=white}},%
  % keep only rows whose column #1 equals #2
  ultra/keep if/.style 2 args={x filter/.append code={%
    \edef\ultratempa{\thisrow{#1}}\edef\ultratempb{#2}%
    \ifx\ultratempa\ultratempb\else\def\pgfmathresult{inf}\fi}},%
}%
%
\begin{figure}[t]
\centering
\begin{tikzpicture}
\begin{axis}[ultra/base, width=0.97\columnwidth, height=4.6cm, xmin=0, xmax=140, ymin=5, ymax=9,
  xlabel={Cumulative training audio (thousand hours)}, ylabel={Full-DEV pooled WER (\%)},
  legend cell align=left, legend columns=2, legend style={at={(.5,-.3)}, anchor=north, column sep=4pt}]
\fill[gray!10] (axis cs:\TrainStageOneEndKhRaw,5) rectangle (axis cs:140,9);
\addplot[ultra/dense, mark options={fill=white}, ultra/keep if={series}{s1c4}]
  table[col sep=comma, x=audio_kh, y=wer] {\paperdatadir training_curve.csv};
\addlegendentry{Stage 1 live $c{=}4$}
\addplot[ultra/dense, dashed, mark=o, ultra/keep if={series}{s1c16}]
  table[col sep=comma, x=audio_kh, y=wer] {\paperdatadir training_curve.csv};
\addlegendentry{Stage 1 $c{=}16$}
\addplot[ultra/rzero, ultra/keep if={series}{s2live}]
  table[col sep=comma, x=audio_kh, y=wer] {\paperdatadir training_curve.csv};
\addlegendentry{Stage 2 live $c{=}4$ (INT8)}
\addplot[ultra/mr, ultra/keep if={series}{s2final}]
  table[col sep=comma, x=audio_kh, y=wer] {\paperdatadir training_curve.csv};
\addlegendentry{Stage 2 final (INT8)}
\node[font=\scriptsize, anchor=north west, okGrey] at (axis cs:\TrainStageOneEndKhRaw,9) {+ Libri-Light};
\end{axis}
\end{tikzpicture}
\caption{\textbf{Training dynamics of the release lineage} on full LibriSpeech DEV (\ultra{} \PriorVer{}).
Stage~1 trains on LibriSpeech 960\,h alone: live WER stops improving after about \TrainPlateauMidKh\,h
(\TrainPlateauMid\% there, \TrainPlateauEnd\% after \TrainStageOneEpochs{} epochs). Stage~2 (shaded) adds
R0-pseudo-labelled Libri-Light and an EMA cooldown, lowering live WER by \TrainLiveGain\pp and final WER by
\TrainFinalGain\pp; the gain arrives with the Libri-Light \texttt{medium} split. Stage-1 points use the PyTorch
streaming decoder, stage-2 points the INT8 engine, which matches it at the stage boundary
(\TrainEngineEnd\% vs.\ \TrainPlateauEnd\%).}
\label{fig:training}
\end{figure}
